Write $\mathfrak m=(t)$, where $t\ne0$ is not a unit. For each $0\ne a\in A$, the intersection hypothesis gives a largest integer $n\geq0$ such that $a\in(t^n)$. Thus $a=t^nu$ with $u\notin\mathfrak m$, so $u$ is a unit. This exponent is unique, since $t$ is not a unit and $A$ is a domain.

Define $v(a)=n$ and $v(a/b)=v(a)-v(b)$ for nonzero $a,b\in A$. Multiplication adds the exponents, so this is well-defined on $\operatorname{Frac}(A)^*$. Also $v(a+b)\geq\min(v(a),v(b))$, since the smaller power of $t$ divides both summands; the same inequality holds for fractions after taking a common denominator. With $v(0)=+\infty$, this defines a valuation. It is onto $\mathbb Z$ because $v(t)=1$, and a nonzero fraction belongs to $A$ exactly when its valuation is nonnegative. Hence $A$ is a discrete valuation ring.
